\documentclass[10pt,a4paper]{amsart}
\usepackage[margin=19mm]{geometry}
\usepackage{amsmath,amssymb,mathtools}
\usepackage[scr=rsfs,cal=euler]{mathalfa}
\usepackage{xfrac}
\usepackage[unicode,hidelinks,bookmarks=false]{hyperref}
\newcommand{\ie}{i.e.\ }
\newcommand{\cf}{cf.\ }
\newcommand{\resp}{respectively\ }
\newcommand{\Subsection}{\subsection}
\providecommand{\nobreakdash}{\nobreak\hbox{-}}
\setlength{\emergencystretch}{2em}
\multlinegap0pt
\usepackage{amssymb}
\usepackage{tikz}
\newtheorem{prop}{Proposition}
\title{Rook characters as symmetric functions}
\author{Rosa Orellana, Alexander N. Wilson, Mike Zabrocki}
\date{Source: arXiv:2609.09042v1 (CC BY 4.0)}
\begin{document}
\maketitle

\noindent\emph{Source and licence.} Rook characters as symmetric functions. Version arXiv:2609.09042v1 (CC BY 4.0), distributed by arXiv under the Creative Commons Attribution 4.0 International licence, http://creativecommons.org/licenses/by/4.0/, as recorded in the arXiv licence metadata for this exact version and shown on the abstract page https://arxiv.org/abs/2609.09042v1. Full licence notice: \texttt{licenses/arxiv-2609.09042-CC-BY-4.0.txt}.\par\smallskip

\noindent\textit{Editorial scope.} Counted unit: the article's prop prop:pleth\_irred\_structure (source0.tex lines 1138--1144). The article's complete original proof of it is delivered with it (source0.tex lines 1146--1160, one \texttt{proof} environment of 15 lines). No mathematics is written, repaired or re-derived: the statement and the proof are the article's own lines, in the article's own notation, and no retained line is substituted or re-flowed.  Retained uncounted as the source-local context the proof needs: lines 951--954; lines 969--971.  Nothing was omitted from any retained span, so the package carries no omission record.  No retained line contains a citation, so the wrapper carries no bibliography and no bibliography entry.  No retained line contains a figure environment, an image-inclusion command or drawing code, so no image is delivered and the package declares no asset dependency.  Editorial formatting only, in addition: an AMS \texttt{amsart} preamble that restates the article's own declarations and macros used by the retained lines, namely the environments \texttt{prop} on separate counters, together with the standard packages those lines need.  The retained environments print in this document as Proposition 1 in order of appearance, whereas the article prints the same items with the numbering of its own full document.  The complete native source of the article as distributed in the arXiv e-print for this exact version is bundled unchanged as \texttt{sources/209772/source0.tex}.
\begin{equation}\label{eq:Schur_coproducts}
c^\lambda_{\nu\mu} = \left< s_\nu[X]s_\mu[Y], s_\lambda[X+Y] \right>~~~~~\hbox{and}~~~~~
g_{\lambda\mu\nu} = \left< s_\nu[X]s_\mu[Y], s_\lambda[XY] \right>~.
\end{equation}

\begin{equation}\label{eq:Cauchy_properties}
\Omega[X+Y] = \Omega[X]\Omega[Y]\qquad\hbox{and}\qquad\Omega[XY] = \sum_{\lambda} s_\lambda[X] s_\lambda[Y]~.
\end{equation}

\begin{prop}
\label{prop:pleth_irred_structure}
For partitions $\alpha, \beta, \gamma$ such that $0 \leq |\gamma| \leqslant |\alpha|+|\beta|$,
\[
\overline{g}_{\alpha\beta}^\gamma = \left< s_\alpha[X] s_\beta[Y], s_\gamma[X+Y+XY] \Omega[XY] \right>
\]
\end{prop}

\begin{proof} We use the identities \eqref{eq:Schur_coproducts}
and \eqref{eq:Cauchy_properties}
where in our sums below each is implicitly over the partitions that are free.
\begin{align*}
&\left< s_\alpha[X] s_\beta[Y], s_\gamma[X+Y+XY] \Omega[XY] \right>\\
&=\sum c^\gamma_{\zeta\rho\sigma} \left< s_\alpha[X] s_\beta[Y], s_\sigma[X]s_\rho[Y]s_\zeta[XY] \Omega[XY] \right>\\
&= \sum c^\gamma_{\zeta\rho\sigma}
g_{\delta\epsilon\zeta} \left< s_\alpha[X] s_\beta[Y], s_\sigma[X]s_\rho[Y]s_\delta[X]s_\epsilon[Y] \Omega[XY] \right>\\
&= \sum c^\gamma_{\zeta\rho\sigma}
g_{\delta\epsilon\zeta} \left< s_\alpha[X] s_\beta[Y], s_\sigma[X]s_\rho[Y]s_\delta[X]s_\epsilon[Y] s_\tau[X] s_\tau[Y] \right>\\
&=\sum
g_{\delta\epsilon\zeta} c^\alpha_{\delta\sigma\tau} c^\beta_{\epsilon\rho\tau} c^\gamma_{\zeta\rho\sigma}
= \overline{g}_{\alpha\beta}^\gamma~. \qedhere
\end{align*}
\end{proof}

\end{document}
